\begin{center}
    \section*{\fontsize{20}{20}\selectfont Chapter 1}
\end{center}
\vspace{10mm}
\section{Introduction}
\subsection{Introduction}

Dhaka loses between \textbf{200 and 300 lives to fire every year}, according to the annual statistics published by the Bangladesh Fire Service and Civil Defence. The loss figures climb sharply in the years that include a single large incident, such as the \textbf{Nimtali fire of 2010}, the \textbf{Tazreen Fashions factory fire of 2012}, the \textbf{Chawkbazar chemical-store fire of 2019}, and the \textbf{FR Tower office blaze of 2019} \cite{fscdAnnual2024}. In each case the fire had been burning for ten minutes or longer before the first call reached a fire station, and in each case the responders who arrived first had no live information about the seat of the fire, the presence of trapped occupants, or the storage of flammable goods in the building. The problem is not a shortage of cameras. Bangladesh's cities are saturated with closed-circuit television, and most factories and commercial buildings are required by law to install smoke detectors. The problem is the \textbf{gap between the device that first sees the fire and the dispatcher who can send a crew}.

This thesis describes an integrated framework that closes that gap. The framework has four layers. The first is an \textbf{autonomous threat-detection layer} that runs a YOLOv8n model exported to ONNX inside a web browser, with a color-space and connected-component heuristic as fallback when no model is available. The second is a \textbf{network-independent alert relay} that pairs a SIP-over-UDP voice call with a peer-to-peer mesh of OpenWrt routers, so that an alert can leave the building even when the public internet and the cellular network are both down. The third is an \textbf{acoustic-wave fire-suppression layer} that drives a low-frequency sine-wave generator in the 30 to 60 Hz band, a regime in which the sound field destabilises the combustion reaction and can extinguish small flames without water or chemicals. The fourth is a \textbf{coordinated emergency-dispatch workflow} that mirrors the operational practice of the Bangladesh Fire Service and Civil Defence: incident intake, automatic station assignment by district, vehicle and crew tracking, escalation to neighbouring stations and to partner agencies, and a citizen-facing channel for reports and feedback. The four layers share a single incident schema and exchange events through a REST API and a server-sent event stream.

The work extends our earlier Amanot-Net project \cite{amanotnet2026}, which targeted weapon detection in unstructured public spaces. Amanot-Net established the edge-first, dual-channel pattern: detection runs on the device, alerts travel over SIP when the network is up and over a low-power wide-area backup when it is not, and the whole system is designed against the assumption that the network will fail at the worst possible moment. The current thesis applies that pattern to fire and smoke, but it goes one step further. Amanot-Net stopped at the alert: a phone rang, a person answered. Fire response does not stop at the alert. The alert has to enter a dispatch workflow, has to be matched to the nearest available crew, has to be tracked until the scene is contained, and the fire itself, when it is small enough, can be attacked directly by an on-site suppressor while the crew is en route. The framework in this thesis therefore adds the suppression layer and the dispatch layer on top of the Amanot-Net pattern, and it treats both as first-class components rather than as external systems that someone else happens to run.

\subsection{Problem Statement}

Fire response in Bangladesh, as in most low- and middle-income countries, runs on three fragile assumptions: that \textbf{someone will see the fire and call the emergency number}, that \textbf{the telephone network will carry the call to the dispatcher}, and that \textbf{the dispatcher will know which station to send}. Each of these assumptions fails often enough that the response time, measured from ignition to first appliance on scene, regularly exceeds fifteen minutes in the older and denser parts of Dhaka \cite{fscdAnnual2024,islam2020dhakafire}. The same period in a Western capital with a comparable population density is closer to five minutes. The gap is not explained by traffic alone.

The \textbf{detection problem} is the first to fail. Most buildings in the informal sector rely on battery-powered smoke detectors that produce a local alarm only. A detector that screams in an empty warehouse at 2 a.m. alerts no one. Where closed-circuit television is installed, it is monitored by a single guard watching many feeds, or it is recorded for forensic review after the fact. The first minute of a fire, when a single extinguisher can still contain it, passes without an external notification in the great majority of cases \cite{shidik2019,sreenu2019}.

The \textbf{relay problem} fails next. The public internet in Bangladesh is good on average but patchy under load. The cellular networks saturate during large public events, during religious festivals, and during the very emergencies in which fire alarms must get through. Cyclones and monsoon floods take cellular base stations offline for hours at a time across whole districts \cite{gsma2023disaster}. A detection system that depends on a working uplink has a single point of failure, and the failure mode is correlated with the failure mode of the event it is supposed to report.

The \textbf{suppression problem} is the one that current academic literature treats the least. Even when detection is fast and the alert is delivered, the fire keeps growing until a crew arrives with water. In a closed room, an electrical panel, a server rack, or a kitchen fryer, water cannot be used, and the chemical extinguishers that can be used require a human to be present. An autonomous suppressor that can act in the first seconds after detection, while the alert is still in flight, can change the size of the fire that the crew eventually faces. \textbf{Acoustic fire suppression}, in which a low-frequency sound field destabilises the flame base, is one of the few non-chemical, non-water suppression mechanisms that can be triggered automatically and that leaves no residue.

The \textbf{dispatch problem} is the slowest to fail and the hardest to fix. The Bangladesh Fire Service and Civil Defence runs a national hotline (the 999 service) and a fleet of approximately 570 stations. Calls arriving at 999 are answered by a telecommunicator who selects a station, calls the station by radio or mobile, and dispatches the crew manually. There is no automatic station assignment, no live vehicle tracking, no shared incident record between the dispatcher and the arriving crew, and no structured handoff to police, ambulance, or utility services when the incident requires more than one agency. The result is that an incident which could be closed by a single engine in eight minutes often stretches to forty, while the dispatcher tries to find a free crew and to convey the address over a noisy phone line.

This thesis addresses all four problems in a single integrated framework. It does not assume that any one of them can be solved in isolation, because the value of the system lies in the chain from detection through suppression to dispatched crew, not in any single link.

\subsection{Problem Background}

The problem background is most easily understood by tracing the four threads that the framework brings together: computer-vision-based fire detection, delay-tolerant and mesh-based alerting, acoustic fire suppression, and computer-aided dispatch in the fire-service context. Each thread has its own research lineage, and the lines cross only rarely.

The \textbf{fire-detection thread} begins with rule-based pixel methods that flagged bright red-yellow regions in a camera frame \cite{celik2010fire}. Convolutional neural networks overtook these methods in accuracy from roughly 2018 onward, with models such as Faster R-CNN, SSD, and the YOLO family demonstrating reliable flame and smoke detection on benchmark datasets like Foggia's and the D-Fire set \cite{dunnhofer2023dfire}. The current state of the art for resource-constrained edge deployment is the \textbf{YOLOv8n variant}, which reaches 0.85 to 0.93 mAP@0.5 on standard fire-smoke benchmarks at 30 to 60 frames per second on a Raspberry Pi 4 \cite{ultralytics2023yolov8,yasar2024fireyolo}. The detection thread has therefore matured to the point where the remaining work is in deployment, not in model design.

The \textbf{alerting thread} has a longer history. Delay-tolerant networking was formalized by Vahdat and Becker in 2000, and the Bundle Protocol was standardized at the IETF in 2007 and revised in 2022 \cite{vahdat2000,rfc9171}. Epidemic routing, spray-and-wait, and encounter-based forwarding were the principal academic contributions of the following decade \cite{spyropoulos2005}. The Serval project demonstrated practical mesh telephony on smartphones, and Briar demonstrated peer-to-peer messaging over Tor hidden services \cite{serval2013,briar2018}. The Mesh-Chat project \cite{meshchat2026} showed that an offline peer-to-peer messenger can run on a stock OpenWrt router with a 19-kilobyte package and a 1.7-megabyte resident memory footprint. These projects collectively establish that a router-based, network-independent relay is feasible; none of them have been wired into a fire-detection pipeline.

The \textbf{acoustic suppression thread} is the youngest of the four. The interaction between sound fields and flames was first reported by \textbf{Lord Rayleigh} in the 19th century, but the modern literature begins with the 2012 demonstration by DARPA and the University of Virginia, in which a pair of large loudspeakers emitting tones in the 30 to 60 Hz band extinguished a small flame at close range \cite{darpa2012acoustic}. Subsequent work by Kim, Park, and others \cite{kim2015acoustic, mckinney2019sonic} established that the dominant mechanism is \textbf{velocity coupling at the flame base}: the oscillating flow thins the flame front beyond the extinction thickness, while the periodic pressure variation disrupts the heat-feedback loop between the flame and the fuel surface. The method is most effective on small, contained flames (pan fires, liquid pool fires, candle-class flames) and less effective on large, ventilated fires. To our knowledge, no prior work has integrated an acoustic suppressor into an autonomous fire-detection pipeline.

The \textbf{dispatch thread} is the oldest of the four. Computer-aided dispatch systems have been in continuous use in North American public-safety answering points since the 1980s, and the Next Generation 911 program in the United States has spent more than a decade migrating those systems to IP-based, geospatially aware platforms \cite{ng9112023}. Equivalent modernization in South Asia is partial. India's 112 emergency number consolidates police, fire, and ambulance calls at state-level command centres, but vehicle tracking and inter-agency escalation remain works in progress \cite{india1122022}. Bangladesh's 999 service, launched in 2017, handles the call-taking well but stops short of automatic dispatch, live tracking, or structured inter-agency coordination \cite{bangladesh9992018}.

The intersection of these four threads is almost empty. There are fire-detection systems that push notifications to a phone app over the internet. There are delay-tolerant messaging systems that work off-grid. There are laboratory demonstrations of acoustic flame extinguishment. There are computer-aided dispatch systems that manage fire response. There is no system, to our knowledge, that puts a YOLOv8-based detector on the edge, a SIP-plus-mesh relay in the middle, an acoustic suppressor at the source, and a structured dispatch workflow at the back end, all sharing a single incident schema. The framework in this thesis is a step toward filling that intersection.

\subsection{Research Objectives}

The overall objective of this research is to design, build, and evaluate an integrated framework that detects fire in its first minute, attempts local suppression while the alert is in flight, relays the alert even when the public network is down, and feeds the alert into a dispatch workflow that assigns a crew and tracks the incident to closure. The specific objectives are listed below.

\begin{itemize}
\item \textbf{Detection.} Design and train a YOLOv8n-based fire and smoke detector that runs at interactive frame rates in a web browser through onnxruntime-web, with a heuristic fallback that produces a usable result when no model file is present.
\item \textbf{Relay.} Build a network-independent alert relay that places a SIP-over-UDP voice call to the on-call responder and, in parallel, floods a compact JSON alert across a peer-to-peer mesh of OpenWrt routers, so that the alert reaches the dispatcher over at least one channel when either the public internet or the cellular network is unavailable.
\item \textbf{Suppression.} Build an acoustic-wave fire suppressor that, on a confirmed detection, emits a low-frequency sine wave in the 30 to 60 Hz band through a connected loudspeaker, with the goal of extinguishing small flames before the crew arrives while the alert is in flight.
\item \textbf{Dispatch.} Design an emergency-dispatch workflow that accepts alerts from the detection and relay layers, assigns each incident to the nearest available fire station by district and geo-distance, tracks the dispatched vehicles and crew, supports escalation to neighbouring stations and to partner agencies, and exposes the resulting incident record through a web portal for citizens, officers, and command-centre operators.
\item \textbf{Integration.} Define a single incident schema that is shared across the four layers, so that an event originating at a camera in a factory in Gazipur can arrive at the dispatch portal in Dhaka with its detection confidence, GPS coordinates, threat type, suppression status, and timestamp intact.
\item \textbf{Evaluation.} Evaluate the framework end to end on detection accuracy, suppression effectiveness on a small alcohol pan fire, alert-delivery latency under both live and severed uplink, dispatch triage time under load, and mesh-relay bandwidth consumption and delivery ratio at scale.
\end{itemize}

\subsection{Motivations}

The most direct motivation for this work is operational. We spent four months in 2025 building Amanot-Net, a weapon-detection system for public spaces. When that project closed, we asked the Bangladesh Fire Service and Civil Defence's Dhaka headquarters whether the same pattern would help them. Their answer, in essence, was that the detection was useful but the alert was useless if it did not enter their dispatch workflow, and that a fire that grew for ten minutes while the crew was en route was a fire that had already won. A phone call to a dispatcher who is already on another call, or who does not know which station to send, buys very little. The framework in this thesis is the result of taking that feedback seriously and extending the Amanot-Net pattern with a suppression layer and a dispatch layer.

A second motivation comes from the observation that the public internet in Bangladesh fails in patterns, not at random. It fails during the hours after a large fire, when everyone in the neighbourhood tries to call their family at once. It fails during cyclones, when the power grid drops and the cellular towers run out of diesel. It fails during political unrest, when the regulator shuts down mobile data on purpose. A fire-response system that cannot tolerate these failures is a system that does not work when it is most needed. The mesh relay in this thesis is a direct response to that observation.

A third motivation is the existence of a suppression mechanism that is fast, clean, and automatable. Water ruins electronics. Dry chemicals ruin kitchens. Inert gas suppression requires sealed enclosures and large storage cylinders. \textbf{Acoustic suppression}, by contrast, needs only a loudspeaker and a few seconds of power, leaves no residue, and can be triggered the instant a detection is confirmed. It cannot extinguish a room fire, but it can extinguish the small flame that would have become a room fire. That window, between the first detection and the arrival of the crew, is the window the acoustic suppressor is designed to fill.

A fourth motivation is cost. Industrial fire-detection systems are priced for export to markets that can absorb a five-figure dollar installation cost per building. The great majority of buildings in Bangladesh's informal sector cannot. A framework that runs its detection layer on a phone or a Raspberry Pi, its relay layer on a fifty-dollar OpenWrt router that the building already owns, its suppression layer on a hundred-dollar public-address loudspeaker, and its dispatch layer on a free Next.js instance, can be deployed at a price point that matches the local economy. The motivation is therefore not only technical but economic.

\subsection{Flow of the Research}

The research followed the four-stage flow shown in Figure~\ref{fig:flow}. The first stage was a literature review covering fire-detection models, delay-tolerant networking, acoustic fire suppression, and computer-aided dispatch, with a particular focus on Bangladesh-specific operational reports. The second stage was the design and training of the YOLOv8n detector on a 73{,}600-image fire and smoke subset of the BD-Threats-v1 corpus, and the parallel development of the heuristic fallback. The third stage was the implementation of the SIP-over-UDP voice call, the OpenWrt mesh relay, and the acoustic-wave suppressor, including the gossip-based flooding protocol, the digest-authentication logic for SIP, and the audio-oscillator tuning for the 30 to 60 Hz band. The fourth stage was the implementation of the dispatch portal and the end-to-end evaluation, in which alerts generated by the detection layer triggered the suppressor and were sent through the relay into the portal, and the portal's response was measured under both live and severed uplink conditions.

\begin{figure}[H]
\centering
\includegraphics[width=\textwidth]{figures/fig5_research_flow.png}
\caption{The four-stage research flow. Stage 1 (M1--M2): literature review and problem definition. Stage 2 (M3--M4): dataset collection and annotation of the BD-Threats-v1 fire subset. Stage 3 (M5--M6): YOLOv8n model training and optimisation, with the HSV heuristic fallback and the acoustic-wave suppressor developed in parallel. Stage 4 (M7--M10): relay and dispatch module development, system integration, and end-to-end evaluation on a 30-node mesh testbed and a small-flame suppression rig.}
\label{fig:flow}
\end{figure}

\subsection{Significance of the Research}

The research is significant on four fronts. The \textbf{first is technical}. The framework is, to our knowledge, the first to wire a YOLOv8-based edge detector to an acoustic-wave suppressor, a SIP-and-mesh relay, and a structured dispatch workflow in a single pipeline with a shared incident schema. The technical contribution is not the invention of any one of these components, all of which build on prior work, but the integration and the schema. The \textbf{second front is operational}. The framework gives the Bangladesh Fire Service and Civil Defence a working prototype of an end-to-end pipeline that they can deploy in pilot form without purchasing industrial-grade equipment, and it gives them a schema they can extend as their modernization program continues. The \textbf{third front is methodological}. The work demonstrates that the edge-first, dual-channel pattern established in Amanot-Net for security applications transfers cleanly to the fire-response domain, and it argues that the pattern is general enough to apply to flood, building-collapse, and road-accident response in future work. The \textbf{fourth front is in the suppression layer itself}. By adding an on-site acoustic suppressor to the pipeline, the framework opens a research direction that the fire-detection literature has so far ignored: the fire can be attacked in the seconds before the crew arrives, not merely observed.

\subsection{Research Contribution}

This thesis makes the following contributions.

\begin{itemize}
\item An integrated framework that combines a YOLOv8n-based fire and smoke detector running in the browser, an acoustic-wave fire suppressor, a SIP-and-mesh alert relay, and a structured dispatch workflow in a single pipeline. The framework is implemented and deployed as two web applications, FireGuard and FireServiceBD, that share a single incident schema.
\item A network-independent relay design in which a SIP-over-UDP voice call is placed in parallel with a peer-to-peer mesh flood of compact JSON alerts. The mesh relay is built on OpenWrt and reuses the Mesh-Chat codebase, extended with Ed25519 message signing and gossip-based flooding.
\item An acoustic-wave suppression design in which a Web Audio oscillator drives a low-frequency sine wave in the 30 to 60 Hz band, validated on a 15-cm alcohol pan fire with a mean extinction time of 4.6 seconds at 45 Hz.
\item A dispatch workflow that mirrors the operational practice of the Bangladesh Fire Service and Civil Defence, including district-based station assignment, multi-agency escalation, and a citizen-facing SOS channel. The workflow is implemented as a Next.js web portal with a Prisma-managed PostgreSQL schema.
\item An evaluation of the framework on a 73{,}600-image fire and smoke subset of BD-Threats-v1, on a 30-node mesh testbed, on a small-flame suppression rig, and on a dispatch portal under simulated concurrent load. The evaluation reports mAP, suppression time, latency, bandwidth, delivery ratio, and triage time.
\item A documented incident schema and REST API that other groups can reuse to build compatible detection, suppression, relay, or dispatch components.
\end{itemize}

\subsection{Complex Engineering Analysis}

The framework sits at the intersection of computer vision, embedded systems, telephony, mesh networking, acoustics, and web engineering. Each layer has its own constraints and its own failure modes, and the layers interact in non-obvious ways. This subsection summarises the most significant engineering problems and the activities undertaken to address them.

\subsubsection{Complex Engineering Problems}

\textbf{Accuracy-latency trade-off.} The detection layer has to run at interactive frame rates on a Raspberry Pi 4 or in a mid-range phone browser, while still producing a useful bounding box. YOLOv8s reaches higher mAP but its inference time on the Pi 4 is above 90 milliseconds, which caps the loop at 11 frames per second and starves the heuristic of CPU. YOLOv8n trades roughly four points of mAP for a 3$\times$ speed-up, which we judged to be the right exchange for a system whose downstream consumer is a dispatch portal and a suppressor trigger rather than a forensic log.

\textbf{Network failure correlation.} The cellular network and the public internet fail together during the events that produce the most fire incidents. A backup channel that depends on either of them is therefore not a backup. The mesh relay addresses this by using the local Wi-Fi spectrum, which is independent of both, but the mesh itself depends on the routers having power. The framework therefore assumes battery-backed routers at the relay layer, which is itself a deployment constraint.

\textbf{Suppression-flame matching.} The acoustic suppressor is effective on small, contained flames (pan fires, candle-class flames, liquid pool fires up to roughly 20 cm) and ineffective on large ventilated fires. The framework therefore has to decide, on the basis of the detection confidence and the bounding-box size, whether to fire the suppressor or to skip straight to the alert. A false suppression is harmless; a missed suppression is no worse than the no-suppressor baseline. The decision rule we adopted is to fire the suppressor whenever the bounding-box area is below 15 percent of the frame, which corresponds to a flame small enough to be in the suppressor's effective regime.

\textbf{Lack of standardised design rules.} There is no international standard for an end-to-end fire-response pipeline that spans detection, suppression, relay, and dispatch. The framework's incident schema is therefore a synthesis drawn from the Next Generation 911 data model, the ETSI emergency communications specifications, and the Bangladesh Fire Service's own internal forms. It is not a standard, but it is documented.

\textbf{Multi-domain system integration.} The detection, suppression, relay, and dispatch layers are written in different languages (TypeScript, Lua, TypeScript) and run on different platforms (browser, OpenWrt, Node.js). The integration challenge is to ensure that an event originating in the browser can trigger the suppressor, traverse the relay, and arrive at the dispatch portal with no field lost or altered. This is solved by a single canonical JSON incident object, validated at every boundary.

\textbf{Privacy, cost, and usability.} The framework processes camera frames on the device and discards them after inference, which addresses the privacy concern at the source. It runs on commodity hardware, which addresses cost. It exposes the dispatch workflow through a web portal that requires no install, which addresses usability. The trade-off is that the system cannot produce forensic video evidence after the fact, which is a deliberate choice.

\subsubsection{Complex Engineering Activities}

\textbf{Dataset development.} We curated a 73{,}600-image fire and smoke subset of the BD-Threats-v1 corpus, drawn from the D-Fire public set, the Foggia benchmark, and original photography in Dhaka. The subset was re-annotated to two classes (fire, smoke) and split 70/20/10 into train, validation, and test.

\textbf{Model training and optimization.} YOLOv8n was trained with transfer learning from the COCO weights, with input resolution 320$\times$320 to match the on-device inference budget. We ran four ablations on input resolution, anchor design, and augmentation, and selected the configuration that gave the best mAP@0.5 on the validation split at sub-40-millisecond inference time on the Pi 4.

\textbf{Suppressor rig construction.} We built a small-flame suppression test rig consisting of a 15-cm-diameter alcohol pan, a 50-watt full-range loudspeaker driven by a Web Audio oscillator through a 100-watt class-D amplifier, and a high-speed camera for extinction-time measurement. The rig was used to characterise the suppressor's extinction time as a function of frequency, intensity, wave type, and flame size.

\textbf{System integration.} The detector, the suppressor, the SIP caller, the mesh relay, and the dispatch portal were integrated through a REST API and a server-sent event stream. The integration was tested with an end-to-end harness that drives a synthetic camera feed, fires the suppressor, observes the resulting SIP call and mesh flood, and asserts that the dispatch portal receives the incident within a configurable latency bound.

\textbf{Failure-scenario testing.} The uplink was intentionally severed at the relay layer in a controlled test, and the framework was required to deliver the alert through the mesh within five seconds. The test was run with 10, 20, and 30 mesh nodes, with and without node churn, and with 0, 10, and 20 percent random packet loss.

\textbf{Innovation and social impact.} The framework's rule-based validation step, which suppresses alerts that do not persist for at least two consecutive frames above the confidence threshold, was designed to reduce the false-alarm rate that has historically made automated fire-detection systems unpopular with responders. The acoustic suppressor extends the system's role from observer to actor, which is a step that the fire-detection literature has so far been reluctant to take. The system's design also considers the privacy of building occupants, the cost of deployment, and the training burden on dispatchers.

\subsection{Thesis Organization}

This thesis is organised into six chapters.

\textbf{Chapter 1} introduces the research problem, its background, the objectives and motivations behind the study, its significance and contributions, and the complex engineering analysis underpinning the framework.

\textbf{Chapter 2} reviews related work in fire and smoke detection, acoustic fire suppression, delay-tolerant and mesh-based messaging, and computer-aided dispatch for fire services, and identifies the gaps in existing research that this thesis addresses.

\textbf{Chapter 3} describes the framework in detail: the detection layer, the suppression layer, the relay layer, the dispatch layer, the shared incident schema, and the algorithms that bind them.

\textbf{Chapter 4} presents the experimental setup and the performance evaluation, including detection accuracy, suppression effectiveness, alert-delivery latency, mesh-relay bandwidth and delivery ratio, dispatch triage time, and an ablation study validating key design choices.

\textbf{Chapter 5} discusses sustainability standards, ethical considerations, and the practical challenges and constraints affecting the system, along with the project timeline.

\textbf{Chapter 6} concludes the thesis, summarising its contributions and outcomes, acknowledging current limitations, and outlining directions for future work.

\subsection{Summary}

This chapter has set out the need for an integrated framework that detects fire in its first minute, attempts local suppression while the alert is in flight, relays the alert when the public network is down, and feeds the alert into a structured dispatch workflow. The next chapter reviews the literature on which the framework is built.
